% TOP_PROBLEM: {"id": 20003353, "title": "Strict special-point sequences, Galois component weights, and the correct Haar target", "category_id": 20, "category": {"id": 20, "name": "miscellaneous", "display_name": "Miscellaneous", "description": "Problems whose source classification does not fit the main mathematical categories.", "slug": "miscellaneous", "order_index": 20, "created_at": "2026-07-31T15:26:25.671Z"}, "status": "partially_solved"}
% FIRST_POSTED: 2026-09-25
% ENTRY_KIND: statement_only
% !TeX program = lualatex
% Definition and sources only; this file is not an LLM solution attempt.
\documentclass[11pt]{article}
\usepackage[margin=25mm]{geometry}
\usepackage{fontspec}
\setmainfont{Latin Modern Roman}
\usepackage{amsmath,amssymb,mathtools}
\usepackage{unicode-math}
\setmathfont{Latin Modern Math}
\usepackage{xurl,hyperref}
\hypersetup{colorlinks=true,urlcolor=blue}
\setcounter{secnumdepth}{0}
\setlength{\parindent}{0pt}
\setlength{\parskip}{5pt}
\setlength{\emergencystretch}{3em}
\begin{document}
\section*{\texorpdfstring{Strict special-point sequences, Galois component weights, and the correct Haar target}{Strict special-point sequences, Galois component weights, and the correct Haar target}}\label{20003353}
\subsection{Definitions and mathematical statement}
A special point on a Shimura variety is a zero-dimensional special subvariety, equivalently a point represented by a Shimura datum with toric Mumford--Tate group. In the canonical-model and connected-component convention of the source, let $\mathcal O_j$ be the finite Galois orbit of $x_j$ and set
$\mu_j=|\mathcal O_j|^{-1}\sum_{x\in\mathcal O_j}\delta_x$. Assume that each fixed proper special subvariety contains only finitely many $x_j$. The assertion is weak convergence
\[
 \mu_j\Longrightarrow\mu_{\mathrm{Haar}},
\]
where the ambient invariant measure is normalized to total mass one. The condition is not that a single finite exceptional set of indices excludes all special subvarieties simultaneously. The field of Galois action and component measure must follow the cited setup, especially for a disconnected ambient Shimura variety.
\par\smallskip\textit{Formulation and concept references:} \href{https://aimath.org/WWN/measrigid/measrigid.pdf}{[S1]}.

\subsection{Short English statement}
Do Galois orbits of special points become uniformly distributed when the sequence eventually avoids each fixed proper special subvariety?
\subsection{Sources}
\begin{itemize}
\item[S1]AIM Workshop Participants, Open Problems from the Workshop Emerging Applications of Measure Rigidity, Conjecture 52, 2004.
\url{https://aimath.org/WWN/measrigid/measrigid.pdf}.
\textit{Formulation source.}
\end{itemize}
\end{document}
